% Paper da 2AV (EL1, CESAR School). Template IEEE conference; IEEEtran.cls v1.8b do CTAN.
% Números vêm de results/ (M6); enquanto o sweep não termina, ficam como TODO: pendente.
\documentclass[conference]{IEEEtran}
\usepackage[T1]{fontenc}
\usepackage[utf8]{inputenc}
\usepackage[brazilian]{babel}
\usepackage{amsmath}
\usepackage{graphicx}
\usepackage{xcolor}

\renewcommand{\IEEEkeywordsname}{Palavras-chave}
\newcommand{\todo}[1]{\textcolor{red}{TODO: #1}}

\begin{document}

\title{CNN ou Transformer como \textit{encoder} da U-Net? Comparação pareada entre ResNet-34 e
MiT-B0 na segmentação de pólipos no Kvasir-SEG}

\author{\IEEEauthorblockN{Lucas Rosati}
\IEEEauthorblockA{CESAR School\\
Recife, PE, Brasil}}

\maketitle

\begin{abstract}
A segmentação de pólipos em imagens de colonoscopia apoia a triagem do câncer colorretal.
Este trabalho compara dois \textit{encoders} pré-treinados no ImageNet para uma U-Net: a CNN
ResNet-34 e o Transformer hierárquico MiT-B0. Os dois cenários são treinados com o mesmo
protocolo no Kvasir-SEG (1000 imagens, divisão fixa 800/100/100) em 30 execuções pareadas por
semente (60 treinos), e o Dice no conjunto de teste é comparado pelo teste \textit{t} pareado
bilateral e pelo teste de Wilcoxon pareado ($\alpha = 0{,}05$). O \textit{encoder} ResNet-34
obteve Dice de \todo{pendente} e o MiT-B0, de \todo{pendente}; a diferença média foi de
\todo{pendente} (IC 95\%: \todo{pendente}; $p = $ \todo{pendente}), \todo{pendente: decisão
sobre H0}. Também são relatados o número de parâmetros e a latência de inferência de cada
cenário.
\end{abstract}

\begin{IEEEkeywords}
segmentação de pólipos, U-Net, \textit{transfer learning}, Transformers, Kvasir-SEG, teste
\textit{t} pareado
\end{IEEEkeywords}

\end{document}
